\begin{reviewcomment}{4}
	The trusted proof service is still an assumption of trust.
	The privacy contribution should be viewed in the context that the TPS has access to confidentiality-sensitive information and builds the witness.
	The implementation also co-locates the TPS with the provider, concentrating this trust.
	Thus, a compromised or colluding TPS may replace attributes or refuse to compute a proof.
	This problem is acknowledged in the discussion, but the practical implications require a more comprehensive discussion.
	The authors should clarify the difference between the cryptographic guarantees provided by TPS-ZKPI and the guarantees based on honest TPS behavior.
	A small threat/guarantee matrix would help clarify things.
\end{reviewcomment}

\begin{reviewresponse}{4}
	We agree, and we have made the distinction explicit in two places.

	First, the trust-model subsection now separates the two kinds of guarantee in the text.
	Authorization correctness is enforced by the commitment and the on-chain proof verification: a proof that does not satisfy the committed relation is rejected regardless of how the TPS is instantiated, and, following the revisions for Comments~1--3, the relation now binds the program length, the authorization context, and the issuer-signed profile claims, so a compromised TPS cannot shorten the program, replay a proof under fresh parameters, or alter signed claims.
	Policy confidentiality, in contrast, is not cryptographically enforced against the TPS: the TPS holds the plaintext policy material and the requester evidence, and a compromised or colluding TPS can expose them, substitute unsigned evidence, or refuse to prove.
	The provider-operated TPS used in the evaluation concentrates this trust at the provider.

	Second, we added the requested matrix (the threat/guarantee table in the revised Problem and Threat Model section).
	Its rows are the threats---passive public-chain observer, untrusted requester, honest-but-curious TPS, compromised or colluding TPS, provider-operated TPS, and forged or unauthorized issuer---and its columns the four guarantees the workflow provides: policy confidentiality, authorization correctness, liveness, and accountability.
	Each cell states whether the guarantee is cryptographic (C), rests on the TPS following its workflow (T), is an external or lifecycle function (E), or is not provided (--).
	The matrix makes the practical implications visible at a glance: against a compromised TPS the prototype keeps authorization correctness for the committed relation but loses confidentiality, liveness (a refusal is indistinguishable from a legitimate denial), and accountability; the Discussion lists accountable proof obligations and reduced-trust proving (collaborative SNARKs, attested execution) as the mechanisms that would recover them.

	\RevisionLocation{Section ``Problem and Threat Model,'' subsection ``System and Trust Model,'' and the new threat/guarantee table.}
	\begin{revisiontext}
		Two guarantees must be distinguished here.
		Authorization correctness is primarily enforced by the commitment and on-chain proof verification: a proof that does not satisfy the committed policy relation is rejected, largely independently of how the TPS is instantiated.
		Policy confidentiality, in contrast, is not cryptographically enforced against the TPS itself; it rests entirely on the trust assumption stated above.
		[\ldots]
		The threat/guarantee table summarizes, for each party or failure mode, which of the four guarantees---policy confidentiality, authorization correctness, liveness, and accountability---is enforced cryptographically, which holds only under the TPS trust assumption, and which the prototype leaves to external functions or does not provide.
	\end{revisiontext}
\end{reviewresponse}
